\documentclass[../../General.tex]{subfiles}
\begin{document}

\section{Evaluación técnica}


\begin{table}[H]
	\captionsetup{justification=raggedright, singlelinecheck=false, labelsep=none}
	
	\caption[Matriz de Evaluación Técnica de Conceptos de Solución actualizada.]{\textit{Matriz de Evaluación Técnica de Conceptos de Solución actualizada.}} \par\medskip
	
	\centering
	\footnotesize
	\renewcommand{\arraystretch}{1.3}
	\begin{tabular}{|p{4.2cm}|c|c|c|c|c|c|c|c|c|}
		\hline
		\rowcolor[HTML]{D9E1F2} 
		& & \multicolumn{2}{c|}{\textbf{Solución 1}} & \multicolumn{2}{c|}{\textbf{Solución 2}} & \multicolumn{2}{c|}{\textbf{Solución 3}} & \multicolumn{2}{c|}{\textbf{Sol. Ideal}} \\ \cline{3-10}
		\rowcolor[HTML]{D9E1F2} 
		\multirow{-2}*{\textbf{Criterio de Evaluación}} & \multirow{-2}*{\textbf{Peso (G)}} & \textbf{P} & \textbf{GP} & \textbf{P} & \textbf{GP} & \textbf{P} & \textbf{GP} & \textbf{P} & \textbf{GP} \\ \hline
		\textbf{Geometría:} (E) & 8 & 3 & 24 & 4 & 32 & 2 & 16 & 4 & 32 \\ \hline
		\textbf{Energía:} (E) & 9 & 2 & 18 & 4 & 36 & 2 & 18 & 4 & 36 \\ \hline
		\textbf{Materia y Térmica:} (E) & 10 & 1 & 10 & 3 & 30 & 4 & 40 & 4 & 40 \\ \hline
		\textbf{Señales:} (E) & 10 & 2 & 20 & 3 & 30 & 4 & 40 & 4 & 40 \\ \hline
		\textbf{Control:} (E) & 9 & 2 & 18 & 3 & 27 & 4 & 36 & 4 & 36 \\ \hline
		\textbf{Comunicación:} (E) & 8 & 2 & 16 & 4 & 32 & 3 & 24 & 4 & 32 \\ \hline
		\textbf{Facilidad de mant.:} (E) & 8 & 2 & 16 & 2 & 16 & 4 & 32 & 4 & 32 \\ \hline\hline
		
		\multicolumn{2}{|r|}{\textbf{Puntaje Total ($Pt$)}} & \multicolumn{2}{c|}{\textbf{122}} & \multicolumn{2}{c|}{\textbf{203}} & \multicolumn{2}{c|}{\textbf{206}} & \multicolumn{2}{c|}{\textbf{248}} \\ \hline
		\multicolumn{2}{|r|}{\textbf{Valor Técnico ($X_i$)}} & \multicolumn{2}{c|}{\textbf{0.49}} & \multicolumn{2}{c|}{\textbf{0.82}} & \multicolumn{2}{c|}{\textbf{0.83}} & \multicolumn{2}{c|}{\textbf{1.00}} \\ \hline
	\end{tabular} \par\medskip
	
	\raggedright
	\scriptsize \textit{Nota.} Elaboración propia. G = Peso del requerimiento; P = Puntaje (escala del 0 al 4); GP = Puntaje Ponderado ($P \times G$); E = Exigencia; D = Deseo. El Valor Técnico ($X_i$) se calcula dividiendo el Puntaje Total de la solución entre el Puntaje Total de la Solución Ideal ($Pt / 248$).
	\label{tab:evaluacion_tecnica}
\end{table}
%
%%

\begin{enumerate}
	\item \textbf{Geometría:} Este criterio evalúa directamente el grado de cumplimiento del requerimiento número 1 de la Lista de Requerimientos (Tabla  \ref{tab:lista_requerimientos}). 
	Para evaluar el volumen y masa del equipo, la solución 2 destaca con un puntaje de 4, ya que su chasis de aluminio \textit{V-Slot} garantiza un diseño compacto y ligero. La solución 1 obtiene 3, pues aunque usa MDF, requiere mayores espesores para rigidez. La solución 3 obtiene 2, al correr el riesgo de acercarse al límite de masa por la alta densidad del acero.
	\item \textbf{Energía:} Este criterio evalúa directamente el grado de cumplimiento del requerimiento número 5 de la Lista de Requerimientos.
	La solución 2, con puntaje de 4, ofrece la alternativa con mayor densidad de energía, mientras que la solución 1 y 3, con puntaje de 2, penalizan la autonomía al tener menor eficiencia energética.
	\item \textbf{Materia:} Este criterio evalúa directamente el grado de cumplimiento del requerimiento número 6 de la Lista de Requerimientos.
	La solución 3, con puntaje de 4, posee el chasis de acero con mayor protección frente al exterior. La solución 2, de puntaje 3, es robusta pero vulnerable al polvo en sus juntas, y la solución 1, de puntaje 1, es deficiente frente a condiciones adversas de temperatura y humedad.
	\item \textbf{Señales:} Este criterio evalúa directamente el grado de cumplimiento del requerimiento número 7 de la Lista de Requerimientos.
	En cuanto a la precisión exigida por la norma IEC 61724-1, la solución 3 es netamente superior, con puntaje 4, logrando integrar instrumentos y controladores de grado industrial. La solución 2 obtiene un buen puntaje de 3, ya que posee una resolución destacable, pero depende de un controlador y protocolos de comunicación con menor robustez. La solución 1 resulta menos eficiente, con un puntaje de 2, pues sus alternativas cuentas con menos fiabilidad en control y almacenamiento de datos.
	\item \textbf{Control:}  Este criterio evalúa directamente el grado de cumplimiento del requerimiento número 8 de la Lista de Requerimientos.
	Para evaluar la capacidad de evadir obstáculos y recalcular rutas de forma dinámica en laderas, la solución 3 domina con un puntaje de 4, apoyándose en un computador industrial y un PLC que procesan nubes de puntos espaciales en tiempo real. La solución 2, con puntaje 3, cuenta con una placa SBC de buen rendimiento, pero su algoritmo de flujo óptico pierde gran parte de su efectividad ante el polvo ambiental de la montaña. A la solución 1 se le otorga 2 puntos al ser la más limitada computacionalmente, con un algoritmo basado en umbrales que puede generar falsos positivos.
	
	\item \textbf{Comunicación:} Este criterio evalúa directamente el grado de cumplimiento del requerimiento número 9 de la Lista de Requerimientos.
	Al evaluar la estabilidad y transmisión de la telemetría, la solución 2 es la alternativa más atractiva, con un puntaje de 4, al incorporar un módulo \textit{LoRa}, excelente para alcanzar grandes distancias en parques extensos consumiendo un mínimo de batería. La solución 3 recibe un puntaje de 3, pues a pesar de que emplea una conexión celular con mayor ancho de banda, depende enteramente de la cobertura de operadoras locales en zonas remotas. Por su parte, a la solución 1 se le asigna 2 puntos, resultando menos viable porque su red \textit{Wi-Fi} requiere de múltiples nodos repetidores para mantener cobertura a lo largo de toda la planta solar.
	
	\item \textbf{Facilidad de Mantenimiento:} Este criterio evalúa directamente el grado de cumplimiento del requerimiento número 10 de la Lista de Requerimientos, junto con la facilidad y frecuencia para realizar mantenimiento a cada concepto de solución.
	Al unificar seguridad y mantenibilidad, la solución 3, con puntaje 4, es superior al tener como alternativas componentes industriales estándar, que soportan gran cantidad de ciclos de trabajo. La solución 2 y 1, con puntajes de 2, presentan componentes con menor grado de robustez y con mayor riesgo de deterioro durante las operaciones de servicio.
\end{enumerate}


\section{Evaluación económica}

\begin{table}[H]
	\captionsetup{justification=raggedright, singlelinecheck=false, labelsep=none}
	
	\caption[Matriz de Evaluación Económica de Conceptos de Solución.]{\textit{Matriz de Evaluación Económica de Conceptos de Solución.}} \par\medskip
	
	\centering
	\footnotesize
	\renewcommand{\arraystretch}{1.3}
	\begin{tabular}{|p{4.5cm}|c|c|c|c|c|c|c|c|c|}
		\hline
		\rowcolor[HTML]{D9E1F2} 
		& & \multicolumn{2}{c|}{\textbf{Solución 1}} & \multicolumn{2}{c|}{\textbf{Solución 2}} & \multicolumn{2}{c|}{\textbf{Solución 3}} & \multicolumn{2}{c|}{\textbf{Sol. Ideal}} \\ \cline{3-10}
		\rowcolor[HTML]{D9E1F2} 
		\multirow{-2}*{\textbf{Criterio de Evaluación}} & \multirow{-2}*{\textbf{Peso (G)}} & \textbf{P} & \textbf{GP} & \textbf{P} & \textbf{GP} & \textbf{P} & \textbf{GP} & \textbf{P} & \textbf{GP} \\ \hline
		\textbf{Costo de Tecnología} & 4 & 3 & 12 & 3 & 12 & 2 & 8 & 4 & 16 \\ \hline
		\textbf{Costo de Mantenimiento} & 3 & 4 & 12 & 2 & 6 & 2 & 6 & 4 & 12 \\ \hline
		\textbf{Disponibilidad de Materiales} & 4 & 4 & 16 & 3 & 12 & 1 & 4 & 4 & 16 \\ \hline
		\textbf{Costo de Fabricación} & 3 & 3 & 9 & 2 & 6 & 1 & 3 & 4 & 12 \\ \hline\hline
		
		\multicolumn{1}{|r|}{\textbf{Puntaje Total ($Pt$)}} & \textbf{14} & \multicolumn{2}{c|}{\textbf{49}} & \multicolumn{2}{c|}{\textbf{36}} & \multicolumn{2}{c|}{\textbf{21}} & \multicolumn{2}{c|}{\textbf{56}} \\ \hline
		\multicolumn{2}{|r|}{\textbf{Valor Económico ($Y_i$)}} & \multicolumn{2}{c|}{\textbf{0.88}} & \multicolumn{2}{c|}{\textbf{0.64}} & \multicolumn{2}{c|}{\textbf{0.38}} & \multicolumn{2}{c|}{\textbf{1.00}} \\ \hline
	\end{tabular} \par\medskip
	
	\raggedright
	\scriptsize \textit{Nota.} Elaboración propia. G = Peso del requerimiento; P = Puntaje (escala del 0 al 4); GP = Puntaje Ponderado ($P \times G$). El Valor Económico ($Y_i$) se calcula dividiendo el Puntaje Total de cada alternativa entre el total de la Solución Ideal ($56$).
	\label{tab:matriz_evaluacion_economica}
\end{table}

\begin{enumerate}
	\item \textbf{Costo de Tecnología:} En cuanto a la inversión necesaria para adquirir los componentes, la solución 1 y la solución 2 obtienen un buen puntaje igual a 3, dado que basan su arquitectura en componentes accesibles y módulos electrónicos de costo estandarizado. Por el contrario, a la solución 3 se le asigna un puntaje igual a 2, puesto que la integración de componentes de grado altamente industrial, como un la PLC, incrementa de manera drástica el costo de capital inicial del sistema, reduciendo su viabilidad para una producción en serie.
	
	\item \textbf{Costo de Mantenimiento:} Para evaluar los gastos operativos a largo plazo, la solución 1 destaca con el mejor puntaje igual a 4, debido a que el reemplazo de sus piezas y sensores básicos resulta económico frente a posibles daños por la intemperie. Tanto la solución 2 como la solución 3 reciben un puntaje igual a 2, ya que implican mayores gastos operativos continuos; la solución 2 exigirá rutinas de limpieza frecuentes para la lente de su cámara CMOS, mientras que la solución 3 demandará un mantenimiento mecánico mucho más costoso para la lubricación de sus ruedas de tracción de oruga y la calibración de la instrumentación industrial.
	
	\item \textbf{Disponibilidad de Materiales:} En cuanto a la facilidad para adquirir los repuestos en el mercado, la solución 1 obtiene un puntaje igual a 4, puesto que aprovecha la amplia disponibilidad local de polímeros genéricos, láminas de MDF y ecosistemas electrónicos de hardware libre, evitando procesos de importación. Por su parte, la solución 2 obtiene un puntaje de 3, ya que cuenta con tecnologías más especializadas que no siempre cuentan con un stock local. De forma similar, la solución 3 tiene un puntaje de 1, dado que cuenta con componentes de grado industrial que están sujetas a demoras por la cadena de suministro internacional y requieren procesos complejos de importación.
	
	\item \textbf{Costo de Fabricación:} Finalmente, respecto a la manufactura y ensamblaje, la solución 1 lidera con un puntaje igual a 3, ya que la fabricación basada en corte láser e impresión 3D permite producir piezas únicas con un costo muy bajo, eliminando por completo la necesidad de invertir en herramientas de conformado especializadas. La solución 2 recibe un puntaje de 2; aunque utilice perfiles modulares V-Slot que facilitan el ensamblaje, requiere de procesos de mecanizado para su caja de aluminio, mientras que la solución 3 recibe un puntaje de 1, pues exige una inversión considerable  para los procesos de soldadura y conformado de acero, así como de sus sensores y controladores industriales.
\end{enumerate}

\end{document}